\chapter{Pressure Systems and Compressed Air Risk Assessment}
\label{chap:Pressure Systems and Compressed Air Risk Assessment}

%---------------------------- SECTION ----------------------------
\section{Why this assessment matters in construction}

\paragraph{Pressure systems and compressed air equipment are ubiquitous on construction sites. Compressors power pneumatic tools, drive nail guns and chipping hammers, inflate tyres, and supply breathing air to operatives working in RPE. Pressure vessels --- including receiver tanks on static and portable compressors, propane and butane cylinders, acetylene cylinders used in hot works, hydraulic accumulators in plant and machinery, and pressurised water jetting equipment --- are present across virtually every project type, from domestic extensions to major civil engineering schemes. The failure modes of pressurised equipment, while statistically less common than those of more visible site hazards, are characterised by their explosive energy release: a compressed air receiver operating at 7 bar and failing catastrophically can discharge stored energy equivalent to a significant explosive charge, projecting metallic fragments at velocities capable of causing fatal injury at distances far beyond any conventional exclusion zone.}

\paragraph{The Pressure Systems Safety Regulations 2000 (PSSR 2000) establish the principal regulatory framework for pressure systems in UK workplaces, and their implementation on construction sites is frequently poor. Investigations by the HSE have consistently identified failures in: the written scheme of examination required by PSSR 2000; the examination of pressure vessels by a competent person at the frequencies specified in the scheme; the maintenance of pressure relief valves and pressure gauges; and the management of hired-in plant for which the scheme of examination is the hiring company's responsibility. The consequence of these failures is not merely a regulatory compliance failure --- it is the creation of a catastrophic failure risk that is statistically concentrated in poorly maintained old equipment operated by personnel who have never been briefed on its hazards or the controls required.}

\paragraph{Acetylene deserves particular attention in any construction pressure systems assessment. Unlike other flammable gases stored as a liquid under pressure, acetylene is stored dissolved in acetone within a porous matrix inside the cylinder, and its failure behaviour is uniquely hazardous. When an acetylene cylinder is involved in a fire, the decomposition reaction is self-sustaining and cannot be stopped by cooling alone: a cylinder involved in fire must be treated as a device with an unpredictable and potentially very long delay before catastrophic explosion, and the standard protocol --- cool with water continuously for at least 20 minutes, then carefully move to a safe outdoor location if the cylinder is cool to the touch --- must be embedded in every site emergency procedure where acetylene is used. The HSE's published guidance on acetylene incidents makes clear that firefighters and construction workers have been killed attempting to manage acetylene cylinders without understanding this behaviour.}

\paragraph{Beyond acetylene, the compressed air circuit itself presents hazards that are routinely underestimated. High-pressure air injected beneath the skin through a hose nozzle or coupling failure can cause subcutaneous emphysema, air embolism, and potentially fatal internal injuries. Air pressure directed at the face can cause serious eye and ear injury. The practice of using compressed air to blow dust off clothing --- prohibited in virtually every written safe system of work and repeatedly identified as a significant causal factor in compressed air injury incidents --- persists on construction sites. Pneumatic tools subject to coupling failure under pressure can project couplings at velocities sufficient to penetrate the human body. The risk assessment for pressure systems must address not only the catastrophic failure modes of vessels and cylinders, but also the acute injury hazards of the compressed air circuit and its associated tooling.}

\barquote{``A pressure vessel is not a passive container --- it is a store of mechanical energy that will be released in its entirety, instantaneously, if the vessel fails. The scheme of examination is not paperwork: it is the mechanism by which the hidden degradation of that vessel is detected before the energy is released in an uncontrolled way.''}

\example{Compressor Receiver Explosion}{A portable diesel compressor with an integral 150-litre air receiver is in use on a groundworks site supplying a pneumatic breaker. The compressor was manufactured in the 1980s and has been in service with the contractor for fourteen years. The receiver has not been examined under a written scheme of examination since the company purchased it second-hand six years ago. During operation, the receiver fails catastrophically due to internal corrosion that has reduced the wall thickness to below the safe operating margin. The explosion projects the receiver base plate 30 metres across the site, fatally injuring an operative working adjacent to a trench. Investigation reveals that the pressure relief valve on the receiver had been manually wired open at some point to prevent nuisance tripping, rendering it entirely ineffective. The written scheme of examination required by PSSR 2000 had never been prepared; the receiver had not been examined by a competent person (an insurance engineer or equivalent) since manufacture. The absence of both the scheme and the examination is a direct causal factor in the failure. The hiring of the compressor to a sub-contractor would have transferred the written scheme responsibility; direct ownership without scheme preparation was the management failure.}

%---------------------------- SECTION ----------------------------
\section{Legal and professional framework}

\paragraph{The legal framework for pressure systems and compressed air on construction sites draws on the Pressure Systems Safety Regulations 2000 as the primary instrument, supported by PUWER 1998 for work equipment generally, the Dangerous Substances and Explosive Atmospheres Regulations 2002 (DSEAR 2002) for flammable gas cylinders, and the Acetylene Safety (England and Wales and Scotland) Regulations 2014. Key frameworks relevant to this chapter include: Health and Safety at Work etc.\ Act 1974; Management of Health and Safety at Work Regulations 1999; Pressure Systems Safety Regulations 2000; Provision and Use of Work Equipment Regulations 1998; DSEAR 2002; Acetylene Safety (England and Wales and Scotland) Regulations 2014; HSE guidance HSG136 (Workplace Transport Safety); and the safe use of compressed gases guidance in HSG240.}

\begin{itemize}
  \bitem{Pressure Systems Safety Regulations 2000 (PSSR 2000)}{The primary statutory instrument governing pressure systems in UK workplaces. PSSR 2000 requires that a written scheme of examination is prepared for every relevant pressure system by a competent person, that the system is examined in accordance with the scheme, and that operating limits for the system are established and not exceeded. The regulations impose duties on designers, manufacturers, suppliers, and users of pressure systems. For construction purposes, the most significant duties fall on the user (the employer operating the compressor or pressure vessel on site): the user must ensure that the written scheme is prepared, is valid, and that the examinations it requires are carried out at the specified intervals. Pressure vessels not examined within the scheme period must be taken out of service until the examination is completed.}
  \bitem{Provision and Use of Work Equipment Regulations 1998 (PUWER)}{Applies to all pneumatic tools, compressors, and pressure equipment as work equipment. PUWER requires that equipment is suitable for its intended use, is maintained in a safe condition, is inspected at appropriate intervals, and is used only by trained and competent operators. The selection of pneumatic tools must take account of the specific hazards of high-pressure air injection, whiplash from failed hoses, and noise exposure (see also Chapter~\ref{chap:Silica, Wood Dust and Hazardous Dust Risk Assessment} and Chapter~\ref{chap:Hazardous Substances and COSHH Risk Assessment}).}
  \bitem{Dangerous Substances and Explosive Atmospheres Regulations 2002 (DSEAR)}{Apply to the storage and handling of flammable gas cylinders, including propane, butane, and acetylene. DSEAR requires that the risks arising from dangerous substances are assessed and that appropriate controls are implemented. On construction sites, this includes the segregation of gas cylinder stores from ignition sources, the provision of adequate ventilation, the securing of cylinders in the upright position, and the management of empty cylinders. The assessment must address the consequences of a gas leak in the cylinder storage area.}
  \bitem{Acetylene Safety (England and Wales and Scotland) Regulations 2014}{Place specific duties on users of acetylene, including the requirement to follow the specific acetylene emergency procedures in the event of fire or damage to an acetylene cylinder. The regulations reference the technical guidance published by the British Compressed Gases Association (BCGA) and require that personnel working with acetylene receive specific training in its unique hazard properties.}
  \bitem{HSG136 --- Workplace Transport Safety}{Provides guidance relevant to the movement of gas cylinders and portable compressors on site by fork-lift trucks, pallet trucks, and other means, including the prohibition on rolling cylinders horizontally and the requirement for appropriate transportation equipment.}
  \bitem{HSE guidance on compressed air injury prevention}{The HSE has published specific guidance on compressed air injury, covering the prohibition on blowing down persons with compressed air, the requirement for dead-man couplings on high-pressure hoses, and the safe working practices for pneumatic tool use. This guidance is not itself legislation but represents recognised good practice and will be considered in any enforcement investigation.}
\end{itemize}

%---------------------------- SECTION ----------------------------
\section{Conducting the assessment using the five-step method}

\subsection{Step 1 --- Identify Hazards}
\paragraph{Pressure system hazards on construction sites fall into four categories: catastrophic vessel failure; gas cylinder incidents (fire, explosion, or leakage); compressed air acute injury hazards; and noise and vibration from pneumatic tools. Catastrophic vessel failure hazards arise from: internal corrosion of receiver tanks, especially in older or poorly maintained equipment; external corrosion at drain points and at the base of receivers that sit directly on ground; over-pressurisation due to failed or defeated safety valves; fatigue failure of vessels subjected to repeated pressure cycling; and mechanical damage from plant contact or falling objects. Gas cylinder hazards include: leakage from damaged valves or regulators creating a flammable or asphyxiating atmosphere; cylinder involvement in fire creating risk of BLEVE (boiling liquid expanding vapour explosion); projectile risk from cylinder valve failure; acetylene decomposition following heat exposure; and the specific cascade decomposition risk of acetylene following a fire. Compressed air acute injury hazards include: subcutaneous air injection from high-pressure air directed at the skin; eye and ear injury from compressed air directed at the face; hose coupling ejection under pressure; and whiplash injury from a failed hose under pressure. Pneumatic tool noise hazards include: hearing damage from prolonged exposure to pneumatic breaking, drilling, and chipping operations (see Chapter~\ref{chap:Asbestos Risk Assessment} regarding silica dust from associated activities).}

\subsection{Step 2 --- Decide Who or What Is Affected}
\paragraph{Persons at risk from pressure system hazards include the operators of pneumatic tools; personnel working adjacent to compressors and pressure vessels; operatives in the storage and handling path of gas cylinders; any person within the potential fragmentation zone of a failing vessel; and members of the public adjacent to an outdoor compressor. Acetylene and LPG cylinder incidents present risks to all personnel in the vicinity of the cylinder store and to adjacent building occupants where cylinders are stored in or adjacent to buildings. The emergency services attending a fire or explosion incident involving pressure vessels or gas cylinders are themselves at significant risk, and the site emergency procedure must address the information to be communicated to the fire service on attendance.}

\subsection{Step 3 --- Evaluate Likelihood and Consequence}
\paragraph{A catastrophic compressor receiver failure carries a consequence of fatal or multiple fatalities and severe structural damage, and must be assessed as Very High before controls are applied. However, the likelihood of failure in a properly examined and maintained vessel is very low; it is only in the absence of examination under a written scheme that the likelihood rises to a level that makes the residual risk unacceptable. The assessment must distinguish between the controlled risk of a compliant system and the uncontrolled risk of a system that has not been examined: the two conditions are not merely quantitatively different, they represent qualitatively different risk profiles. Compressed air injury hazards carry a lower consequence for most scenarios but the likelihood of occurrence --- driven by the persistent practice of misusing compressed air --- is significant, and the injury outcome can be life-changing.}

\subsection{Step 4 --- Select and Implement Controls}
\paragraph{Controls for pressure system hazards must address both the condition of the equipment and the behaviour of those using it. For pressure vessels: written scheme of examination prepared by a competent person; examination at the frequencies specified in the scheme; safety valves tested and recorded; pressure gauges calibrated and within date; drain valves operated and recorded; and operating limits posted adjacent to the compressor. For gas cylinders: storage in a purpose-made outdoor cylinder cage with adequate ventilation; segregation of full, empty, and different gas types; cylinders secured upright at all times; flashback arrestors fitted on welding equipment; regulators inspected before use; and emergency procedure for fire or damaged cylinder posted at the storage location. For compressed air circuits: dead-man couplings (self-closing under pressure loss) on all hose connections where practicable; hose whip arrestors at all compressor outlets; prohibition on blowing down persons with compressed air communicated and enforced; and pre-use inspection of hoses for damage and coupling condition.}

\subsection{Step 5 --- Record and Review}
\paragraph{The written scheme of examination is the primary record for PSSR 2000 compliance. It must be held on site or immediately accessible, must identify each pressure vessel and its examination requirements, and must record the date and outcome of each examination. Maintenance records for safety valves, pressure gauges, and drain valves must be maintained. Gas cylinder delivery and removal records must identify the cylinders on site by serial number and date. Toolbox talk records should document the briefing of operatives on compressed air safety, the prohibition on blowing down, and the gas cylinder emergency procedure. The risk assessment must be reviewed when new pressure equipment is introduced to the site, when plant is hired in from a new supplier, and following any incident involving pressure systems or gas cylinders.}

\example{Five-Step Application}{A civil engineering contractor is operating a static diesel compressor with a 300-litre receiver to supply pneumatic concrete breaking tools on a highway improvement project. \textbf{Step~1 (Hazards):} Receiver over-pressurisation, corrosion of receiver base; pneumatic hose coupling ejection; worker misuse of compressed air; noise from breaking tools. \textbf{Step~2 (Who Affected):} Tool operators, banksman, public in adjacent carriageway. \textbf{Step~3 (Evaluation):} Receiver failure rated Very High consequence/Low likelihood (with scheme); compressed air misuse rated Medium consequence/Medium likelihood. \textbf{Step~4 (Controls):} Written scheme prepared by insurance engineer, annual examination completed and certificate on file; safety valve tested monthly; daily pre-start inspection checklist for compressor including gauge reading; dead-man couplings on all hose connections; whip checks fitted at compressor outlets; toolbox talk delivered to all operatives prohibiting use of air to blow down clothing; hearing protection zone established. \textbf{Step~5 (Record):} Scheme and examination certificate on file; daily pre-start records retained; TBT attendance sheet signed.}

%---------------------------- SECTION ----------------------------
\section{Applying the hierarchy of controls}

\paragraph{The hierarchy for pressure systems should first seek to eliminate or reduce the inventory of pressurised equipment and gas cylinders on site, before layering engineering and administrative controls over what remains.}

\begin{itemize}
  \bitem{Elimination}{Consider whether the work can be done without compressed air or pressurised gas. Electrically powered tools eliminate the need for an air compressor; battery-powered tools further reduce trailing leads. Propane salamander heaters can sometimes be substituted with electric heating in smaller enclosed spaces. Gas-free welding processes (MIG, MMA with solid wire) reduce the acetylene inventory.}
  \bitem{Substitution}{Replace older, poorly documented pressure vessels with modern equipment that comes with a current written scheme and examination certificate. Hire new plant from reputable hire companies with current PSSR 2000 documentation. Use hydraulic rather than pneumatic equipment for higher-risk applications where hydraulic systems can provide equivalent function.}
  \bitem{Engineering controls}{Written scheme of examination and regular competent-person examination is the primary engineering safeguard for vessel integrity. Pressure safety valves set below the maximum allowable working pressure (MAWP), tested and certified, prevent over-pressurisation. Dead-man couplings, whip checks, and hose restraints manage the consequences of hose failure under pressure. Flame arrestors and flashback arrestors prevent flame travel back into gas cylinders during hot works.}
  \bitem{Administrative controls}{Permit-to-work system for first use of any new or hired pressure vessel, requiring documentary evidence of the written scheme. Daily pre-start inspection checklists for compressors. Toolbox talks on compressed air safety and gas cylinder emergency procedures. Prohibition of use of compressed air for cleaning clothing or blowing down persons, enforced through supervision. Segregation and labelling of gas cylinder stores.}
  \bitem{PPE}{Eye protection and hearing protection for pneumatic tool operators. Gloves for gas cylinder handling. Appropriate PPE for hot works activities associated with gas cylinder use (see Chapter~\ref{chap:Hot Works and Fire Risk Assessment}).}
\end{itemize}

%---------------------------- SECTION ----------------------------
\section{Cadence of assessment and review triggers}

\paragraph{The cadence of pressure systems risk assessment must be linked to the examination schedule in the written scheme and to the operational lifecycle of the equipment on site.}

\begin{itemize}
  \bitem{At introduction of any new or hired pressure equipment}{Before any new compressor, pressure vessel, or gas cylinder manifold system is put into service on site, documentary evidence of PSSR 2000 compliance must be obtained and reviewed. For hired plant, the hire company bears responsibility for the written scheme; this responsibility must be confirmed in writing at the time of hire. For purchased plant, the employer must commission the written scheme before use.}
  \bitem{At the frequency specified in the written scheme}{Pressure vessels must be examined at the intervals specified in the written scheme, and must be taken out of service if the examination interval expires. Monthly testing of pressure relief valves and weekly inspection of safety-critical components (gauges, drain valves, hose connections) must be recorded.}
  \bitem{Following any pressure system incident or near miss}{Any unexpected pressure release, hose failure, coupling ejection, or safety valve actuation must trigger an immediate inspection of the system and formal review of the risk assessment before the system is returned to service. RIDDOR reporting requirements must be considered: dangerous occurrences under RIDDOR Schedule~2 include the failure of a closed vessel or associated pipe-work forming part of a pressure system.}
  \bitem{On significant change in operating conditions}{If the operating pressure or operating temperature of a system is changed, or if the duty of the system changes significantly, the written scheme must be reviewed and revised if necessary. Changes in compressed air supply pressure resulting from modifications to the distribution circuit must be assessed for compatibility with downstream equipment.}
  \bitem{Following discovery of external corrosion or physical damage}{Any discovery of external corrosion, mechanical damage, or deterioration of a pressure vessel, cylinder, or high-pressure hose must result in the equipment being taken out of service immediately and examined by a competent person before being returned to use.}
\end{itemize}

%---------------------------- SECTION ----------------------------
\section{Common failure modes}

\begin{itemize}
  \bitem{No written scheme of examination for owned pressure vessels}{The most common PSSR 2000 compliance failure on construction sites is the absence of a written scheme of examination for compressor receivers, particularly for older equipment that was purchased second-hand. Without the scheme, the examination cadence is undefined, the examination cannot be conducted in a legally compliant manner, and the growing risk from internal corrosion and component degradation goes entirely undetected.}
  \bitem{Safety valves defeated or over-set}{Nuisance actuation of pressure relief valves --- a symptom of over-pressurisation caused by incorrect pressure switch settings or a failing unloading valve --- is frequently addressed by the simple expedient of wiring the valve open, replacing it with a valve set to a higher pressure, or removing it entirely. Each of these responses converts a protected system into an unprotected one. The correct response to a nuisance-tripping safety valve is to investigate and correct the over-pressurisation cause, not to defeat the protection.}
  \bitem{Compressed air used to blow down clothing and skin}{Despite being prohibited in every safe system of work for compressed air, the practice of blowing dust and debris off clothing and skin using compressed air is observed on construction sites with regularity. High-pressure air directed at the skin can inject air subcutaneously or through even small abrasions, causing air embolism and death. Eye injury from compressed air is also well-documented. The risk assessment must prohibit this practice explicitly, communicate the prohibition to all operatives, and provide alternative cleaning methods (brushes, vacuum tools) so the practice does not simply default back to compressed air.}
  \bitem{Acetylene cylinders stored horizontally or subjected to heat}{Acetylene cylinders must be stored and used in the upright position to prevent acetone migration and to ensure the cylinder valve can be closed in an emergency. Acetylene cylinders stored near heat sources, in direct sunlight, or adjacent to hot works operations are at risk of initiating the decomposition reaction. The storage area for acetylene must be in the shade, adequately ventilated, and distant from any potential ignition source.}
  \bitem{Hired plant received without PSSR documentation review}{Plant and equipment hired in from external suppliers may arrive on site without the written scheme or examination certificate being requested or reviewed. Site managers and plant managers must establish a protocol for verifying PSSR 2000 documentation before any hired pressure vessel is put into service.}
  \bitem{No emergency procedure for a gas cylinder fire or damaged acetylene cylinder}{In the absence of a specific, briefed emergency procedure for a gas cylinder fire or a damaged or overheating acetylene cylinder, personnel attending such an incident may attempt inappropriate interventions --- attempting to move a hot acetylene cylinder, or remaining in the vicinity of a cylinder involved in fire. The emergency procedure must be posted at the cylinder store and briefed to all site personnel.}
\end{itemize}

\example{Failure Pattern}{A roofing contractor uses a portable gas torch burning propane for applying felt to a flat roof. The propane cylinder is stored in the flat roof access cupboard at the base of the external ladder, adjacent to the building fabric. At the end of the working day, the operative closes the torch but fails to close the cylinder valve. Overnight, a small propane leak develops at the worn regulator connection. By the following morning, the concentration of propane in the enclosed access cupboard has reached a flammable level. The first operative arriving on site the following morning opens the cupboard, and the ignition of an accumulated LPG pocket by a static discharge causes a flash fire that burns both operatives in the access area. Investigation reveals that: no gas cylinder risk assessment had been carried out; the regulator had not been inspected; no procedure for cylinder closure at end of work was documented; and the cylinder storage location had no ventilation and was adjacent to an ignition source. The incident is entirely consistent with the preventable pattern described in DSEAR guidance.}

%---------------------------- CHAPTER SUMMARY ----------------------------
\newpage
\section{Chapter Summary}

\begin{table}[H]
\centering
\begin{tabularx}{\textwidth}{>{\raggedright\arraybackslash}p{3.2cm}X}
\toprule
\textbf{Assessment Element} & \textbf{Operational Requirement} \\
\midrule
Legal/Professional Basis & Pressure Systems Safety Regulations 2000; PUWER 1998; DSEAR 2002; Acetylene Safety Regulations 2014; Health and Safety at Work etc.\ Act 1974; Management Regulations 1999. \\
Method & Apply the full five-step process and document inherent/residual risk. Prepare and comply with a written scheme of examination for all relevant pressure systems. \\
Controls & Written scheme of examination with periodic competent-person examination; certified safety valves and pressure gauges; dead-man couplings and hose whip checks; prohibition on misuse of compressed air; gas cylinder storage in ventilated cage segregated from ignition sources; emergency procedure for fire or damaged cylinder. \\
Cadence & On introduction of new equipment; at scheme-specified examination intervals; following any incident; on change in operating conditions; on discovery of damage or corrosion. \\
Assurance & PSSR 2000 examination certificates on file; RIDDOR reporting of dangerous occurrences; toolbox talk records; daily pre-start inspection checklists. \\
\bottomrule
\end{tabularx}
\end{table}

\begin{itemize}
  \bitem{Key takeaway 1}{Every relevant pressure vessel on site must be covered by a written scheme of examination prepared by a competent person and must be examined at the intervals specified. Vessels not examined within the scheme period must be taken out of service. This applies to owned and hired plant.}
  \bitem{Key takeaway 2}{Safety valves must never be wired open, over-set, or removed. The correct response to nuisance safety valve actuation is to investigate and correct the underlying cause of over-pressurisation, not to defeat the protection.}
  \bitem{Key takeaway 3}{Compressed air must never be used to blow down clothing or skin. The prohibition must be explicitly stated in the risk assessment, communicated by toolbox talk, and enforced by supervision.}
  \bitem{Key takeaway 4}{Acetylene cylinders must be stored upright, in a ventilated location away from ignition sources, and managed under the specific emergency procedure for fire or heat exposure. A damaged or overheating acetylene cylinder must never be moved until it is confirmed cool to the touch.}
  \bitem{Key takeaway 5}{Documentation for hired pressure plant --- written scheme, examination certificate, and operating limits --- must be obtained and verified before the equipment is put into service on site. The transfer of PSSR 2000 responsibility to the hire company must be confirmed in writing.}
\end{itemize}

\barquote{In Chapter~\ref{chap:Traffic Management and Logistics Risk Assessment}, we turn to the management of vehicle and pedestrian traffic on and around construction sites --- an environment in which the interaction between moving plant and site personnel remains one of the leading causes of fatal and serious injury in the construction industry.}
